% Example showing pullquote.sty used directly in LaTeX, with no
% Pandoc involved. Compile with lualatex or xelatex for full font
% parity and custom font loading.
%
% pullquote.sty must be on TEXINPUTS or in the same directory as this file;
% it lives at the repo root.
\documentclass[12pt, a4paper]{article}
\usepackage[margin=20mm]{geometry}
\usepackage{parskip}
\usepackage{setspace}
\usepackage{titlesec}
\usepackage{microtype}
\usepackage{pullquote}
\usepackage[colorlinks=true, bookmarks=true, bookmarksnumbered=true]{hyperref}

% Globally set inline code to Maroon
\let\oldtexttt\texttt
\renewcommand{\texttt}[1]{\textcolor{Maroon}{\oldtexttt{#1}}}

\setlength{\emergencystretch}{3em}

% Noto Serif/Noto Sans/Fira Mono, matching test/settings/shared.yaml's
% mainfont/sansfont/monofont exactly.
\setmainfont{Noto Serif}
\setsansfont{Noto Sans}
\setmonofont{Fira Mono}
\newfontfamily\pqheadingfont{Noto Serif}[UprightFont={* SemiBold},ItalicFont={* SemiBold Italic}]

\titleformat*{\section}{\Large\pqheadingfont}
\titleformat*{\subsection}{\large\pqheadingfont}

% Hex-code colors from the "Hex Codes" section below, defined once up front.
\definecolor{pqHexText}{HTML}{827397}
\definecolor{pqHexBar}{HTML}{FFAAAA}

\title{Demonstration of the \texttt{pullquote} LaTeX Package}
\author{Nandakumar Chandrasekhar}
\date{\today}

\begin{document}
\maketitle

\setstretch{1.25}

\section*{Introduction}

This document showcases the features of the \texttt{pullquote} LaTeX package. Each example displays the required syntax alongside its rendered output, ensuring visually identical layouts across LaTeX (PDF), Typst (PDF), and HTML formats.

\section*{Environment Invocations}

In accordance with LaTeX conventions, pullquotes must be instantiated exclusively via the \texttt{pullquote} environment.

\begin{pullquote}
"Typography is the craft of endowing human language with a durable visual form."
--- Robert Bringhurst
\end{pullquote}

\section*{Font Sizing}

The package provides a sizing scale to adjust the text relative to the document's base font size. While the filter supports a full 5-step scale (from \texttt{s} to \texttt{2xl}) alongside custom arbitrary units, the larger sizes below are the most practical for pullquotes. If omitted, it naturally inherits the document's current font size.

\begin{pullquote}[size=m]
\textbf{size=m:} Medium (Base document size)
\end{pullquote}

\begin{pullquote}[size=l]
\textbf{size=l:} Large
\end{pullquote}

\begin{pullquote}[size=xl]
\textbf{size=xl:} Extra Large
\end{pullquote}

\begin{pullquote}[size=2xl]
\textbf{size=2xl:} Extra Extra Large
\end{pullquote}

\begin{pullquote}[size=\fontsize{24pt}{28.8pt}\selectfont]
\textbf{24pt:} Custom Arbitrary Dimensions
\end{pullquote}

\section*{Font Weights and Shapes}

The package provides predefined typographic utility attributes.

\subsection*{Font Weights}

\begin{pullquote}[weight=bold]
This paragraph is in bold weight.
\end{pullquote}

\subsection*{Font Styles}

By default, pullquotes perfectly inherit the surrounding document's font style. You can override this using the style utilities.

\begin{pullquote}[style=italic]
This paragraph is explicitly rendered in italics, overriding the inherited body text style.
\end{pullquote}

\section*{Custom Fonts}

Because pullquotes are designed to fit seamlessly into your existing layout, they automatically inherit your document's ambient font settings. However, if you want a specific quote to visually break away from the body text, you can override it using the \texttt{font} attribute. The package natively wraps this in a \texttt{\textbackslash fontspec} call.

\begin{pullquote}[font={Georgia}]
This pullquote overrides the document defaults and explicitly uses the Georgia font.
\end{pullquote}

\section*{Colors}

This section describes how to apply colors to the pullquote text and its decorative left border.

\subsection*{Basic Color Attributes}

\begin{pullquote}[color=CadetBlue, barcolor=Thistle]
This quote uses the SVG named color \textbf{CadetBlue} for text and the pastel \textbf{Thistle} for the border.
\end{pullquote}

\subsection*{Hex Codes}

\begin{pullquote}[color=pqHexText, barcolor=pqHexBar]
This quote uses raw Hex codes: \textbf{\#827397} for text and the shorthand \textbf{\#FAA} (Soft Coral) for the bar.
\end{pullquote}

\subsection*{Color Mixing}

The package supports LaTeX's \texttt{xcolor} percentage-mixing syntax.

\begin{pullquote}[color=Indigo!90!Black, barcolor=Indigo!20]
This quote uses a three-part blend for the text (\textbf{Indigo mixed at 90\% with Black}) and a standard two-part blend with white for the bar.
\end{pullquote}

\section*{Box Padding and Borders}

The \texttt{paddingleft} attribute controls the gap between the bar and the quote text, independent of the bar's own thickness.

\begin{pullquote}[barwidth=0.5em, paddingleft=1.5em]
This quote widens the gap between the bar and the text using \textbf{paddingleft=1.5em}, independent of the 0.5em bar thickness.
\end{pullquote}

The remaining three sides — \texttt{paddingright}, \texttt{paddingtop}, and \texttt{paddingbottom} — can each be tuned independently of \texttt{paddingleft}, giving full control over the box's inner spacing.

\begin{pullquote}[barcolor=SeaGreen, paddingleft=1em, paddingright=1em, paddingtop=1em, paddingbottom=1em]
This quote sets an even 1em of padding on every side, giving the text room to breathe on all edges rather than just next to the bar.
\end{pullquote}

\section*{Block-Level Positioning and Width}

By adjusting \texttt{width} and \texttt{boxalign}, you can set the width of the pullquote and position the entire block to the left, center, or right of the page.

\begin{pullquote}[width=0.45\linewidth, boxalign=right]
This quote takes up 45\% of the page width and is aligned flush to the right margin.
\end{pullquote}

\section*{Inner Text Alignment}

The \texttt{textalign} attribute controls text alignment \emph{within} the box.

\begin{pullquote}[textalign=center, width=0.7\linewidth, boxalign=left]
This quote takes up 70\% of the page and is aligned to the left margin, but the text \emph{inside} the box is centered.
\end{pullquote}

\section*{Interline Spacing (Line-Height)}

Because \texttt{pullquote.sty} delegates semantic spacing translations to Pandoc, standalone usage requires passing the raw baseline multiplier directly to the \texttt{skip} attribute.

\begin{pullquote}[size=xl, skip=1.8]
This quote uses a \texttt{skip} value of 1.8, acting as a double-spacing multiplier to let the text breathe.
\end{pullquote}

Tightening works the same way by using a lower multiplier:

\begin{pullquote}[size=l, skip=1.2]
This quote is deliberately set to 1.2, making it ideal for dense, heavily styled multi-line blocks.
\end{pullquote}

\section*{Global Metadata Configuration}

You can establish project-wide defaults for your pullquotes using \texttt{\textbackslash pullquotesetup}.

\pullquotesetup{
  color=DarkSlateGray,
  barcolor=CadetBlue,
  width=0.8\linewidth,
  size=l,
  boxalign=center
}

\begin{pullquote}
This quote has no inline attributes. It inherits its styling entirely from the setup declaration defined at the top of this document.
\end{pullquote}

% Restore defaults
\pullquotesetup{
  color=, barcolor=pqdefaultbar, width=0.8\linewidth,
  barwidth=4pt, paddingleft=12pt, paddingright=0pt,
  paddingtop=4pt, paddingbottom=4pt, boxalign=center,
  size=, weight=, style=, font=, textalign=, skip=
}

\section*{Combining Multiple Attributes}

Every key composes independently to create highly customized layouts:

\begin{pullquote}[font={Georgia}, style=italic, size=xl, width=0.8\linewidth, boxalign=center, color=DarkSlateGray, barcolor=SteelBlue, barwidth=4pt, paddingleft=1.5em, skip=1.55]
"Typography is the detail and the presentation of a story. It represents the voice of an atmosphere, or historical setting of some kind. It can do a lot of things."

— Cyrus Highsmith
\end{pullquote}

\end{document}
